\begin{abstract}
This review examines how extensions of classical continuum mechanics change thermomechanical wave
propagation, energy transport and interface phenomena. It organizes generalized thermoelastic, multiphysics,
memory-dependent, fractional, gradient and nonlocal models by constitutive mechanism, governing structure,
predicted modes and validation. For the exact classical Eulerian balance, the outward-positive total-energy flux
is $F_i=q_i-\sigma_{ij}v_j+\mathcal{E}v_i$; the reduced $J_i=q_i-\sigma_{ij}v_j$ is used only for the stated
first-order small-amplitude wave balance. Lord--Shulman and Green--Lindsay are single-temperature formulations,
distinct from two-temperature thermoelasticity: the former relaxes heat flux, whereas the latter introduces
temperature-rate dependence. The Eringen-type nonlocal dispersion relation considered here has bounded frequency
at high wavenumber, and quasi-electrostatic piezoelectric models do not contain a propagating electromagnetic
branch.

The review is constructive rather than bibliographic. Each family is presented through the balance laws it
retains, the constitutive statement that defines it, the assumptions under which that statement is admissible, and
the wave modes it predicts; where a closed form exists, it is used as a benchmark against which an independent
implementation is tested, and where no closed form exists the predicted mode count is stated instead of a
number. A five-axis taxonomy (physics retained, constitutive order, length scale, memory, coupling) and a
model-selection guide let a reader place any new model in the family and see which of its predictions are
already fixed by its own assumptions. The review compares model-specific assumptions and identifies questions for
further analytical, numerical and experimental study.
\end{abstract}

\vspace{1.0em}
\noindent\textbf{Keywords:} generalized thermoelasticity; non-Fourier heat conduction; piezothermoelasticity;
magneto-electro-elasticity; nonlocal and gradient elasticity; memory-dependent derivative; fractional
thermoelasticity; wave propagation; interface and surface waves; energy flux; benchmark verification
